\documentclass[10pt,reqno,a4paper]{amsart}

\usepackage[utf8]{inputenc}
\usepackage[english]{babel}
\usepackage{amsmath,amssymb,amsthm,amsfonts}
\usepackage{mathtools}
\usepackage[margin=25mm]{geometry}
\usepackage{microtype}
\usepackage[colorlinks=true,linkcolor=blue,urlcolor=blue]{hyperref}
\usepackage[capitalize]{cleveref}

\numberwithin{equation}{section}

\theoremstyle{plain}
\newtheorem{theorem}{Theorem}[section]
\newtheorem{proposition}[theorem]{Proposition}
\newtheorem{lemma}[theorem]{Lemma}
\newtheorem{corollary}[theorem]{Corollary}
\theoremstyle{definition}
\newtheorem{definition}[theorem]{Definition}
\theoremstyle{remark}
\newtheorem{remark}[theorem]{Remark}

\newcommand{\R}{\mathbb{R}}
\newcommand{\N}{\mathbb{N}}
\newcommand{\C}{\mathbb{C}}
\newcommand{\1}{\mathbf{1}}
\newcommand{\de}{\,\mathrm{d}}
\newcommand{\abs}[1]{\left|#1\right|}
\newcommand{\norm}[1]{\left\lVert#1\right\rVert}
\newcommand{\Fourier}{\mathcal{F}}
\newcommand{\Ucal}{\mathcal{U}}
\newcommand{\PV}{\operatorname{pv}}
\newcommand{\supp}{\operatorname{supp}}

\title[The Blackwell constant]{Identification of the Blackwell constant:\\
the Blackwell constant by Abel-regularized Fourier inversion}
\author{Benny Avelin}
\date{July 29, 2026}

\begin{document}

\begin{abstract}
  We consider the remaining constant-identification step in the formalization
  of Blackwell's renewal theorem.  The symmetric principal-value calculation
  gives only one half of the constant
  \(1/m\).  The missing half is the boundary term selected by the one-sided
  renewal series.  We give a complete Abel-regularized Fourier argument which
  retains the characteristic-function estimates already formalized, avoids
  the improper Dirichlet integral, and yields the full constant.  We then
  derive a uniform local bound, the interval form of Blackwell's theorem, and
  the key renewal theorem for directly Riemann integrable functions.  We also
  explain why the alternative Laplace--Tauberian route is not shorter.
\end{abstract}

\maketitle

\section{Setting and the formalized input}\label{sec:setting}

Let \(\mu\) be a probability measure on \(\R\), and let
\begin{equation}\label{eq:mean}
  m:=\int_\R x\,\mu(\de x)\in(0,\infty).
\end{equation}
We assume that \(\mu\) is nonlattice, has a finite second moment, and satisfies
\begin{equation}\label{eq:left-exponential}
  L(\theta_0):=\int_\R e^{-\theta_0x}\,\mu(\de x)<1
\end{equation}
for some \(\theta_0>0\).  The last condition is stronger than the usual
hypotheses of Blackwell's theorem, but it holds for the tilted Gaussian
increment law used in Chapter~7.  It implies that the renewal measure
\begin{equation}\label{eq:renewal-measure}
  U:=\sum_{n=0}^{\infty}\mu^{*n}
\end{equation}
is finite on every half-line \((-\infty,b]\).

We use the Fourier convention
\begin{align}
  \widehat w(t)
  &:=
  \int_\R e^{-2\pi itx}w(x)\,\de x, \label{eq:fourier-convention}\\
  \chi(t)
  &:=
  \int_\R e^{itx}\,\mu(\de x). \label{eq:characteristic-function}
\end{align}
Then
\begin{equation}\label{eq:fourier-convolution}
  \Fourier\left(y\mapsto\int_\R w(y-z)\,\mu(\de z)\right)(t)
  =
  \widehat w(t)\chi(-2\pi t).
\end{equation}
The finite second moment gives
\begin{equation}\label{eq:characteristic-expansion}
  1-\chi(-2\pi t)
  =
  2\pi i mt+O(t^2),
  \qquad t\to0.
\end{equation}

The Lean development already contains the following consequence.  Suppose
that \(v\) and \(\widehat v\) are integrable, \(\widehat v\) has compact
support, and
\begin{equation}\label{eq:difference-kernel}
  \abs{\widehat v(t)}\leq C\abs t.
\end{equation}
Then the partial renewal convolutions converge pointwise, and their limit
satisfies
\begin{equation}\label{eq:difference-kernel-limit}
  \sum_{n=0}^{\infty}\int_\R v(y-z)\,\mu^{*n}(\de z)
  \longrightarrow 0,
  \qquad y\to\infty.
\end{equation}
This is Blackwell's theorem for difference kernels.  The condition
\(\widehat v(0)=0\), equivalently \(\int v=0\), cancels the pole in
\cref{eq:characteristic-expansion}.

\section{Why the symmetric pole is not sufficient}\label{sec:principal-value}

For a kernel \(w\) with \(\widehat w(0)\neq0\), the formal resolvent is
\begin{equation}\label{eq:formal-resolvent}
  \frac{\widehat w(t)}{1-\chi(-2\pi t)}
  \sim
  \frac{\widehat w(0)}{2\pi i mt}.
\end{equation}
The right-hand side is not Lebesgue integrable near the origin.  If it is
interpreted as a symmetric principal value, then for \(y>0\)
\begin{align}
  \PV\int_{-T}^{T}
  \frac{\widehat w(0)e^{2\pi iyt}}{2\pi i mt}\,\de t
  &=
  \frac{\widehat w(0)}{\pi m}
  \int_0^T\frac{\sin(2\pi yt)}{t}\,\de t \notag\\
  &\longrightarrow
  \frac{\widehat w(0)}{2m}.
  \label{eq:half-constant}
\end{align}
Thus the Dirichlet integral alone gives one half of the desired limit.

The renewal series is one-sided in the convolution index \(n\).  Its Fourier
resolvent therefore has a boundary prescription.  At the level of
distributions, the model singularity is
\begin{equation}\label{eq:boundary-value}
  \frac{1}{m}\1_{(0,\infty)}
  =
  \Fourier^{-1}
  \left(
    \PV\frac{1}{2\pi i mt}
    +
    \frac{1}{2m}\delta_0
  \right).
\end{equation}
The principal-value term in \cref{eq:boundary-value} gives
\(\frac{1}{2m}\operatorname{sgn}(y)\), while the atom gives the missing
constant \(1/(2m)\).  Any direct pole proof must construct both terms.  A
symmetric improper integral, by itself, cannot identify the Blackwell
constant.

\begin{remark}\label{rem:bochner-integral}
  The uncorrected function in \cref{eq:formal-resolvent} is not Bochner
  integrable.  It therefore cannot be inserted directly into the ordinary
  Lebesgue Fourier transform used by the present Lean theorem.  The
  Abel regularization below stays within ordinary integrals until the final
  limit and produces the boundary term in \cref{eq:boundary-value}
  automatically.
\end{remark}

\section{Abel-regularized Fourier inversion}\label{sec:abel}

For \(0<r<1\), define the finite Abel renewal measure
\begin{equation}\label{eq:abel-renewal}
  U_r:=\sum_{n=0}^{\infty}r^n\mu^{*n}.
\end{equation}
It has total mass \((1-r)^{-1}\).  If \(w,\widehat w\in L^1(\R)\), Fourier
inversion and Tonelli's theorem give
\begin{equation}\label{eq:abel-parseval}
  \int_\R w(y-z)\,U_r(\de z)
  =
  \int_\R
  \frac{\widehat w(t)e^{2\pi iyt}}
       {1-r\chi(-2\pi t)}
  \,\de t.
\end{equation}
There is no singularity in this identity.

We compare the denominator with the model
\begin{equation}\label{eq:model-denominator}
  b_r(t):=(1-r)+2\pi i rmt.
\end{equation}
Set
\begin{equation}\label{eq:corrected-resolvent}
  D_r(t)
  :=
  \frac{1}{1-r\chi(-2\pi t)}
  -
  \frac{1}{(1-r)+2\pi i rmt}.
\end{equation}

\begin{lemma}[Uniform pole correction]\label{lem:uniform-pole-correction}
  Let \(0<T<\infty\).  There are \(r_0\in(0,1)\) and \(C<\infty\) such that
  \begin{equation}\label{eq:uniform-correction-bound}
    \abs{D_r(t)}\leq C,
    \qquad
    r_0\leq r<1,\quad \abs t\leq T.
  \end{equation}
  For every \(t\neq0\),
  \begin{equation}\label{eq:correction-pointwise-limit}
    D_r(t)\longrightarrow
    D_1(t)
    :=
    \frac{1}{1-\chi(-2\pi t)}
    -
    \frac{1}{2\pi i mt}.
  \end{equation}
\end{lemma}

\begin{proof}
  Write
  \begin{equation*}
    q(t):=1-\chi(-2\pi t)-2\pi i mt.
  \end{equation*}
  By \cref{eq:characteristic-expansion}, there are \(\delta>0\) and
  \(C_0<\infty\) such that
  \begin{equation}\label{eq:q-quadratic}
    \abs{q(t)}\leq C_0t^2,
    \qquad \abs t\leq\delta.
  \end{equation}
  After decreasing \(\delta\), the imaginary part in
  \cref{eq:characteristic-expansion} and the nonnegative real part of
  \(1-r\chi\) give
  \begin{align}
    \abs{1-r\chi(-2\pi t)}
    &\geq c\bigl((1-r)+\abs t\bigr), \label{eq:true-denominator-lower}\\
    \abs{(1-r)+2\pi i rmt}
    &\geq c\bigl((1-r)+\abs t\bigr) \label{eq:model-denominator-lower}
  \end{align}
  for \(r\geq1/2\) and \(\abs t\leq\delta\), where \(c=c(m)>0\).
  Since the two denominators differ by \(rq(t)\),
  \begin{equation*}
    \abs{D_r(t)}
    \leq
    \frac{rC_0t^2}
    {c^2((1-r)+\abs t)^2}
    \leq \frac{C_0}{c^2}.
  \end{equation*}

  On the compact annulus \(\delta\leq\abs t\leq T\), nonlatticeness implies
  \(\chi(-2\pi t)\neq1\).  Hence \(1-\chi(-2\pi t)\) is bounded away from zero
  there.  Both denominators in \cref{eq:corrected-resolvent} are therefore
  bounded away from zero for \(r\) sufficiently close to one.  This proves
  \cref{eq:uniform-correction-bound}.  The pointwise convergence in
  \cref{eq:correction-pointwise-limit} is immediate.
\end{proof}

\begin{lemma}[The model pole]\label{lem:model-pole}
  Suppose that \(w,\widehat w\in L^1(\R)\) and Fourier inversion holds for
  \(w\).  Then
  \begin{align}
    M_r(y)
    &:=
    \int_\R
    \frac{\widehat w(t)e^{2\pi iyt}}
         {(1-r)+2\pi i rmt}
    \,\de t \notag\\
    &=
    \frac{1}{rm}
    \int_{-\infty}^{y}
    \exp\left(
      -\frac{1-r}{rm}(y-x)
    \right)
    w(x)\,\de x. \label{eq:model-pole-formula}
  \end{align}
  Consequently,
  \begin{equation}\label{eq:model-pole-limit}
    M_r(y)\longrightarrow
    \frac{1}{m}\int_{-\infty}^{y}w(x)\,\de x
    \qquad\text{as }r\uparrow1.
  \end{equation}
\end{lemma}

\begin{proof}
  We use
  \begin{equation*}
    \frac{1}{(1-r)+2\pi i rmt}
    =
    \int_0^\infty
    e^{-(1-r)s}e^{-2\pi i rmts}\,\de s.
  \end{equation*}
  Fubini's theorem and Fourier inversion yield
  \begin{equation*}
    M_r(y)
    =
    \int_0^\infty e^{-(1-r)s}w(y-rms)\,\de s.
  \end{equation*}
  The change of variables \(x=y-rms\) gives
  \cref{eq:model-pole-formula}.  Dominated convergence gives
  \cref{eq:model-pole-limit}.
\end{proof}

We now fix a continuous real-valued kernel \(\kappa\) satisfying
\begin{equation}\label{eq:reference-kernel}
  \kappa\geq0,\qquad
  \kappa,\widehat\kappa\in L^1(\R),\qquad
  \supp\widehat\kappa\subseteq[-T,T].
\end{equation}
The sinc-squared kernel already formalized in
\texttt{RenewalKernel.lean}, after normalization and scaling, has these
properties.

\begin{proposition}[Reference-kernel limit]\label{prop:reference-kernel}
  Under the assumptions of \cref{sec:setting},
  \begin{equation}\label{eq:reference-kernel-limit}
    \sum_{n=0}^{\infty}
    \int_\R\kappa(y-z)\,\mu^{*n}(\de z)
    \longrightarrow
    \frac{1}{m}\int_\R\kappa(x)\,\de x,
    \qquad y\to\infty.
  \end{equation}
  In particular, the series on the left is finite for every \(y\).
\end{proposition}

\begin{proof}
  By \cref{eq:abel-parseval}, we split the Abel convolution as
  \begin{equation}\label{eq:abel-split}
    \int_\R\kappa(y-z)\,U_r(\de z)
    =
    M_r(y)+E_r(y),
  \end{equation}
  where
  \begin{equation*}
    E_r(y):=
    \int_{-T}^{T}
    \widehat\kappa(t)e^{2\pi iyt}D_r(t)\,\de t.
  \end{equation*}
  By \cref{lem:uniform-pole-correction} and dominated convergence,
  \begin{equation}\label{eq:correction-abel-limit}
    E_r(y)\longrightarrow
    E(y)
    :=
    \int_{-T}^{T}
    \widehat\kappa(t)e^{2\pi iyt}D_1(t)\,\de t
  \end{equation}
  for every \(y\).  The integrand defining \(E\) belongs to \(L^1\), and hence
  the Riemann--Lebesgue lemma gives
  \begin{equation}\label{eq:correction-vanishes}
    E(y)\longrightarrow0,
    \qquad y\to\infty.
  \end{equation}
  By \cref{lem:model-pole},
  \begin{equation}\label{eq:model-abel-limit-reference}
    M_r(y)\longrightarrow
    \frac{1}{m}\int_{-\infty}^{y}\kappa(x)\,\de x.
  \end{equation}

  Since \(\kappa\geq0\), monotone convergence in the series index gives
  \begin{equation*}
    \lim_{r\uparrow1}
    \int_\R\kappa(y-z)\,U_r(\de z)
    =
    \sum_{n=0}^{\infty}
    \int_\R\kappa(y-z)\,\mu^{*n}(\de z).
  \end{equation*}
  Combining this identity with
  \cref{eq:abel-split,eq:correction-abel-limit,eq:model-abel-limit-reference}
  shows that
  \begin{equation}\label{eq:reference-representation}
    \sum_{n=0}^{\infty}
    \int_\R\kappa(y-z)\,\mu^{*n}(\de z)
    =
    \frac{1}{m}\int_{-\infty}^{y}\kappa(x)\,\de x+E(y).
  \end{equation}
  We let \(y\to\infty\) and use \cref{eq:correction-vanishes}.
\end{proof}

\begin{remark}[Where the missing half enters]\label{rem:missing-half}
  Formula \cref{eq:model-pole-formula} converges to
  \(m^{-1}\int_{-\infty}^y w\), not to
  \((2m)^{-1}\int w\).  Thus the Abel denominator in
  \cref{eq:model-denominator} contains both the principal-value term and the
  atom in \cref{eq:boundary-value}.  No separate Dirichlet-integral
  calculation is needed.
\end{remark}

\section{General band-limited kernels}\label{sec:band-limited}

Let \(w\) be a continuous real- or complex-valued kernel such that
\(w,\widehat w\in L^1(\R)\), \(\widehat w\) has compact support, and
\(\widehat w\) is Lipschitz at the origin.  Choose a reference kernel
\(\kappa\) as in \cref{eq:reference-kernel}, scaled so that its Fourier support
contains that of \(w\), and set
\begin{equation}\label{eq:difference-decomposition}
  a:=\frac{\int_\R w}{\int_\R\kappa},
  \qquad
  v:=w-a\kappa.
\end{equation}
Then \(\widehat v(0)=0\).  The Lipschitz assumption and compact Fourier support
give
\begin{equation*}
  \abs{\widehat v(t)}\leq C\abs t,
  \qquad t\in\R.
\end{equation*}
The formalized difference-kernel theorem applies to \(v\), while
\cref{prop:reference-kernel} applies to \(\kappa\).  We obtain the following
statement.

\begin{theorem}[Band-limited key renewal limit]\label{thm:band-limited}
  Under the assumptions of \cref{sec:setting},
  \begin{equation}\label{eq:band-limited-limit}
    \sum_{n=0}^{\infty}
    \int_\R w(y-z)\,\mu^{*n}(\de z)
    \longrightarrow
    \frac{1}{m}\int_\R w(x)\,\de x,
    \qquad y\to\infty,
  \end{equation}
  for every kernel \(w\) in the preceding class.
\end{theorem}

\begin{proof}
  For every finite \(N\), the partial renewal convolution of \(w\) is the sum
  of the corresponding convolution of \(v\) and \(a\kappa\).  The first
  sequence converges by the formalized \(N\to\infty\) theorem for difference
  kernels, and its limit tends to zero as \(y\to\infty\).  The second sequence
  converges by monotone convergence and
  \cref{prop:reference-kernel}.  Since
  \(a\int\kappa=\int w\), this gives \cref{eq:band-limited-limit}.
\end{proof}

\section{Uniform local boundedness}\label{sec:local-bound}

We choose the scale of the nonnegative reference kernel so that
\begin{equation}\label{eq:kernel-dominates-interval}
  \kappa(x)\geq c_\kappa>0,
  \qquad \abs x\leq1.
\end{equation}
This is possible because the normalized sinc-squared kernel is continuous and
strictly positive in a neighborhood of the origin.

\begin{proposition}[Uniform local bound]\label{prop:uniform-local-bound}
  There exists \(C_U<\infty\) such that
  \begin{equation}\label{eq:uniform-local-bound}
    U([y-1,y+1])\leq C_U,
    \qquad y\in\R.
  \end{equation}
\end{proposition}

\begin{proof}
  By \cref{eq:kernel-dominates-interval},
  \begin{equation*}
    c_\kappa U([y-1,y+1])
    \leq
    \int_\R\kappa(y-z)\,U(\de z).
  \end{equation*}
  The right-hand side converges as \(y\to\infty\) by
  \cref{prop:reference-kernel}, and is therefore bounded for all sufficiently
  large \(y\).  For \(y\) bounded above, the interval \([y-1,y+1]\) is
  contained in one fixed half-line.  Its \(U\)-mass is finite by
  \cref{eq:left-exponential}.  This proves \cref{eq:uniform-local-bound}.
\end{proof}

For a measurable function \(g\), set
\begin{equation}\label{eq:dri-norm}
  \norm{g}_{\mathrm{dri}}
  :=
  \sum_{j\in\mathbb Z}
  \sup_{x\in[j,j+1]}\abs{g(x)}.
\end{equation}
The local bound gives
\begin{equation}\label{eq:dri-convolution-bound}
  \sup_{y\in\R}
  \int_\R\abs{g(y-z)}\,U(\de z)
  \leq
  C_U\norm{g}_{\mathrm{dri}},
\end{equation}
after enlarging \(C_U\).

\section{Blackwell's theorem and the key renewal theorem}\label{sec:key-renewal}

Let \(\rho\) be a normalized nonnegative sinc-squared kernel, and put
\begin{equation}\label{eq:approximate-identity}
  \rho_a(x):=a\rho(ax),
  \qquad a\geq1.
\end{equation}
Then \(\int\rho_a=1\), the Fourier transform of \(\rho_a\) has compact
support, and \((\rho_a)_{a\geq1}\) is an approximate identity.

\begin{lemma}[Band-limited approximation]\label{lem:band-limited-approximation}
  If \(f\) is continuous and compactly supported, then
  \begin{equation}\label{eq:dri-approximation}
    \norm{f*\rho_a-f}_{\mathrm{dri}}\longrightarrow0,
    \qquad a\to\infty.
  \end{equation}
\end{lemma}

\begin{proof}
  On a fixed compact interval, \cref{eq:dri-approximation} follows from
  uniform continuity and the approximate-identity property.  Outside a
  slightly larger interval, the sinc-squared bound
  \begin{equation*}
    \rho_a(x)\leq\frac{Ca}{1+a^2x^2}
  \end{equation*}
  gives a summable cellwise majorant whose sum tends to zero as
  \(a\to\infty\).  Combining the compact and tail estimates proves the claim.
\end{proof}

\begin{proposition}[Compactly supported test functions]\label{prop:compact-tests}
  If \(f\) is continuous and compactly supported, then
  \begin{equation}\label{eq:compact-test-limit}
    \int_\R f(y-z)\,U(\de z)
    \longrightarrow
    \frac{1}{m}\int_\R f(x)\,\de x,
    \qquad y\to\infty.
  \end{equation}
\end{proposition}

\begin{proof}
  The function \(f*\rho_a\) is band-limited and satisfies the hypotheses of
  \cref{thm:band-limited}.  By \cref{eq:dri-convolution-bound},
  \begin{equation*}
    \sup_y
    \abs{
      \int_\R
      \bigl((f*\rho_a)-f\bigr)(y-z)\,U(\de z)
    }
    \leq
    C_U\norm{f*\rho_a-f}_{\mathrm{dri}}.
  \end{equation*}
  We first let \(y\to\infty\) and then \(a\to\infty\), using
  \cref{lem:band-limited-approximation} and
  \(\int f*\rho_a=\int f\).
\end{proof}

\begin{theorem}[Blackwell]\label{thm:blackwell}
  For every \(h>0\),
  \begin{equation}\label{eq:blackwell}
    U((y,y+h])
    \longrightarrow
    \frac{h}{m},
    \qquad y\to\infty.
  \end{equation}
\end{theorem}

\begin{proof}
  For every \(\varepsilon>0\), choose nonnegative continuous compactly
  supported functions \(f_\varepsilon^-\) and \(f_\varepsilon^+\) such that
  \begin{equation*}
    f_\varepsilon^-
    \leq\1_{[0,h)}
    \leq f_\varepsilon^+,
    \qquad
    \int_\R f_\varepsilon^-\geq h-\varepsilon,
    \qquad
    \int_\R f_\varepsilon^+\leq h+\varepsilon.
  \end{equation*}
  Apply \cref{prop:compact-tests} to both functions and squeeze.  Since
  \begin{equation*}
    \int_\R \1_{[0,h)}(y+h-z)\,U(\de z)=U((y,y+h]),
  \end{equation*}
  letting \(\varepsilon\downarrow0\) gives \cref{eq:blackwell}.
\end{proof}

\begin{definition}\label{def:directly-riemann-integrable}
  A function \(g:\R\to\C\) is directly Riemann integrable if it is almost
  everywhere continuous and
  \begin{equation}\label{eq:dri-condition}
    \sum_{j\in\mathbb Z}
    \sup_{x\in[j,j+1]}\abs{g(x)}<\infty.
  \end{equation}
\end{definition}

\begin{theorem}[Key renewal theorem]\label{thm:key-renewal}
  If \(g\) is directly Riemann integrable, then
  \begin{equation}\label{eq:key-renewal}
    \int_\R g(y-z)\,U(\de z)
    \longrightarrow
    \frac{1}{m}\int_\R g(x)\,\de x,
    \qquad y\to\infty.
  \end{equation}
\end{theorem}

\begin{proof}
  It is enough to consider \(g\geq0\).  Fix a mesh \(\delta>0\), and bound
  \(g\) from below and above by its lower and upper step functions on the
  intervals \([k\delta,(k+1)\delta)\).  For each fixed cell,
  \cref{thm:blackwell} gives
  \begin{equation*}
    U((y-(k+1)\delta,y-k\delta])
    \longrightarrow
    \frac{\delta}{m}.
  \end{equation*}
  The uniform local bound in \cref{prop:uniform-local-bound} and
  \cref{eq:dri-condition} control the tails uniformly in \(y\).  We may
  therefore pass to the limit first on finitely many cells and then in the
  cell cutoff.  The resulting lower and upper bounds are the corresponding
  Riemann sums divided by \(m\).

  On each compact interval, almost-everywhere continuity is the Lebesgue
  criterion for Riemann integrability.  The summability in
  \cref{eq:dri-condition} controls the two tails.  Hence the lower and upper
  sums converge to \(\int g\) as \(\delta\downarrow0\), proving
  \cref{eq:key-renewal}.  For a complex-valued \(g\), apply the result to the
  positive and negative parts of its real and imaginary parts.
\end{proof}

\section{Comparison with the Laplace--Tauberian route}\label{sec:tauberian}

The exponential-transform identities already formalized give
\begin{equation}\label{eq:renewal-laplace}
  \int_\R e^{-\theta x}\,U(\de x)
  =
  \frac{1}{1-L(\theta)}.
\end{equation}
The finite mean implies
\begin{equation}\label{eq:laplace-asymptotic}
  1-L(\theta)\sim m\theta,
  \qquad \theta\downarrow0.
\end{equation}
Since \(U((-\infty,0])<\infty\), a monotone Tauberian theorem would yield
\begin{equation}\label{eq:elementary-renewal}
  U([0,x])\sim\frac{x}{m}.
\end{equation}
There are two remaining difficulties.

First, \cref{eq:elementary-renewal} does not follow from a one-line
exponential squeeze.  Taking \(\theta=a/x\) only gives a bound containing
\(e^a/a\).  The exact constant requires the index-one Karamata theorem, or
its proof by polynomial approximation of indicator functions in the variable
\(e^{-x}\).

Second, the cumulative asymptotic \cref{eq:elementary-renewal} does not imply
the local limit \cref{eq:blackwell}.  One still needs nonlatticeness, uniform
local control, and a smoothing or Tauberian argument at fixed spatial scale.
Using the already formalized difference-kernel theorem at this stage does not
remove the need for local compactness.  In particular, asymptotic invariance
under each fixed translation together with a Cesàro limit is not sufficient
for pointwise convergence.

Thus the Laplace route requires both a new monotone Tauberian lemma and a
second local-limit argument.  It does not consume the pole-correction work
already present in the branch.

\end{document}
